\documentclass[a4paper,11pt]{article}

\usepackage[utf8]{inputenc}
\usepackage[T1]{fontenc}
\usepackage[ngerman]{babel}
\usepackage{amsmath,amssymb,amsthm}
\usepackage{geometry}
\geometry{margin=2.3cm}
\usepackage{enumitem}
\usepackage{fancyhdr}
\usepackage{xcolor}
\usepackage{hyperref}

%================= Dark-Mode Farben (Catppuccin Mocha) =================
\definecolor{base}{HTML}{1E1E2E}
\definecolor{fg}{HTML}{CDD6F4}
\definecolor{accent}{HTML}{89B4FA}   % Überschriften
\definecolor{good}{HTML}{A6E3A1}     % richtig / Merke
\definecolor{bad}{HTML}{F38BA8}      % Fehler / Achtung
\definecolor{rulecol}{HTML}{45475A}
\pagecolor{base}
\makeatletter
\newenvironment{psmallmatrix}{\left(\begin{smallmatrix}}{\end{smallmatrix}\right)}
\makeatother

\hypersetup{colorlinks=true, linkcolor=accent, urlcolor=accent}

\newcommand{\transp}{^{\mathsf{T}}}
\newcommand{\norm}[1]{\left\lVert #1 \right\rVert}
\newcommand{\inner}[2]{\left\langle #1,\, #2 \right\rangle}
\newcommand{\rang}{\operatorname{rang}}
\newcommand{\spur}{\operatorname{spur}}
\newcommand{\diag}{\operatorname{diag}}

%================= Theorem-/Box-Stile =================
\newtheoremstyle{def0}{6pt}{6pt}{}{}{\bfseries\color{accent}}{}{.5em}{}
\theoremstyle{def0}
\newtheorem*{defn}{Definition}
\newtheorem*{merke}{\textcolor{good}{Merke}}
\newtheorem*{fehler}{\textcolor{bad}{Dein Fehler im Blatt}}
\newtheorem*{bsp}{Beispiel}

%================= Kopf-/Fußzeile, Überschriften =================
\usepackage{titlesec}
\titleformat{\section}{\color{accent}\normalfont\Large\bfseries}{\thesection}{0.6em}{}
\titleformat{\subsection}{\color{accent}\normalfont\large\bfseries}{\thesubsection}{0.6em}{}

\pagestyle{fancy}
\fancyhf{}
\lhead{\color{fg}\small Lernzettel -- Matrizen}
\rhead{\color{fg}\small Nach Übungsblatt 01}
\cfoot{\color{fg}\thepage}
\renewcommand{\headrule}{{\color{rulecol}\hrule height 0.4pt}}

\setlength{\parindent}{0pt}
\setlength{\parskip}{0.4em}

\begin{document}
\color{fg}
% hyperref/toc Linkfarbe + Inhaltsverzeichnis hell
\begingroup
\begin{center}
  {\color{accent}\LARGE\bfseries Lernzettel: Deine Schwachstellen}\\[0.3em]
  {\large Fortgeschrittene Lineare Algebra und Analytische Geometrie}\\[0.2em]
  {\small erstellt aus deinem ausgefüllten Übungsblatt 01}
\end{center}
\endgroup

\vspace{0.4em}{\color{rulecol}\hrule}\vspace{0.6em}

Dieser Zettel behandelt \emph{gezielt} die Themen, bei denen auf deinem Blatt etwas
\textbf{konzeptionell} schiefging, plus die zwei Themen, die du nochmal erklärt haben
wolltest (Finite Differenzen, Regression). Jeder Abschnitt: \emph{Intuition} $\to$
\emph{Definition/Formel} $\to$ \emph{durchgerechnetes Beispiel} $\to$ \emph{dein Fehler}.
Alle Formeln stehen auch im Open-Book-Zettel -- die Verweise stehen jeweils dabei.

\tableofcontents
\vspace{0.4em}{\color{rulecol}\hrule}

%==================================================================
\section{Positiv definite Matrizen \& Definitheit}

\textbf{Intuition.} „Positiv definit“ (PD) heißt: die Matrix $A$ „zieht“ jeden Vektor
auf seine eigene Seite -- $\vec{x}\transp A\vec{x} > 0$ für alle $\vec{x}\ne\vec{0}$.
Das ist die Matrix-Verallgemeinerung von „eine positive Zahl“.

\begin{defn}
Eine \emph{symmetrische} Matrix $A$ heißt \textbf{positiv definit}, wenn
$\vec{x}\transp A\vec{x} > 0$ für alle $\vec{x}\ne\vec{0}$.
\end{defn}

\textbf{Zwei Tests (Open-Book S.\,2 „Definitheit“ \& S.\,1 „LU“):}
\begin{itemize}[noitemsep]
  \item \textbf{Eigenwerte:} $A$ ist PD $\iff$ \emph{alle} Eigenwerte $\lambda_i > 0$.
        (PSD: $\lambda_i\ge0$; indefinit: positive \emph{und} negative.)
  \item \textbf{Führende Hauptminoren:} alle führenden Unterdeterminanten
        $\Delta_1,\Delta_2,\dots,\Delta_n > 0$.
\end{itemize}

\begin{bsp}
$A=\begin{psmallmatrix}4&2&2\\2&5&5\\2&5&14\end{psmallmatrix}$:\quad
$\Delta_1=4>0$,\;
$\Delta_2=\det\begin{psmallmatrix}4&2\\2&5\end{psmallmatrix}=16>0$,\;
$\Delta_3=\det A = 144>0$ $\Rightarrow$ PD.
\end{bsp}

\begin{fehler}
In A1c und A6a hast du gesagt, eine positiv definite Matrix \emph{könne auch negative
Eigenwerte haben} bzw.\ du hast Definitheit über „$A\ge0$ für jedes Element“ geprüft.
\textbf{Beides ist falsch:} Definitheit hat \emph{nichts} mit den einzelnen Einträgen zu
tun. Prüfe \textbf{Eigenwerte} ($>0$) oder \textbf{Hauptminoren} ($>0$). PD $\Rightarrow$
\emph{ausschließlich} positive Eigenwerte (A1c ist also \emph{wahr}).
\end{fehler}

%==================================================================
\section{Diagonalisierbarkeit: wann existiert \texorpdfstring{$A=S\Lambda S^{-1}$}{A=SLS-1}?}

\textbf{Intuition.} Diagonalisieren heißt: eine Basis aus Eigenvektoren finden, in der
$A$ nur noch streckt (Diagonalmatrix $\Lambda$). Das geht aber \emph{nur}, wenn man
genug Eigenvektoren hat, um ganz $\mathbb{R}^n$ aufzuspannen.

\begin{merke}
$A\in\mathbb{R}^{n\times n}$ ist diagonalisierbar ($A=S\Lambda S^{-1}$ mit
$S=[\vec v_1\;\dots\;\vec v_n]$) $\iff$ es gibt $n$ \textbf{linear unabhängige}
Eigenvektoren. Sonst existiert die Zerlegung \emph{nicht}.\\
\textbf{Spezialfall:} $A$ symmetrisch $\Rightarrow$ \emph{immer} diagonalisierbar, sogar
orthogonal: $A=Q\Lambda Q\transp$ mit $Q^{-1}=Q\transp$ (Spektralsatz).
\end{merke}

\begin{bsp}
$A=\begin{psmallmatrix}1&1\\0&1\end{psmallmatrix}$ hat nur den Eigenwert $\lambda=1$ mit
\emph{einem} Eigenvektor $(1,0)\transp$ -- es fehlt ein zweiter unabhängiger.
$A$ ist \textbf{nicht} diagonalisierbar.
\end{bsp}

\begin{fehler}
In A1d hast du „Ja“ angekreuzt (jede quadratische Matrix sei zerlegbar). Das ist
\textbf{falsch}: nur diagonalisierbare Matrizen. Merksatz: „brauche $n$ unabhängige
Eigenvektoren“ (Open-Book S.\,1 „$S\Lambda S^{-1}$“, Limit-Zeile).
\end{fehler}

%==================================================================
\section{Untervektorraum vs.\ affiner Raum (Hyperebenen)}

\textbf{Intuition.} Ein Untervektorraum (UVR) muss den \textbf{Ursprung} enthalten -- er
geht „durch null“. Eine verschobene Gerade/Ebene ist zwar flach, aber \emph{kein} UVR.

\begin{merke}
$U$ ist UVR $\iff$ (1) $\vec0\in U$, (2) abgeschlossen unter $+$, (3) abgeschlossen
unter Skalarmultiplikation. \textbf{Schnelltest:} Geht das Gebilde durch den Ursprung?
\end{merke}

Eine Hyperebene $E:\ \vec n\transp\vec x + n_0 = 0$ geht genau dann durch den Ursprung,
wenn $n_0=0$ (denn $\vec n\transp\vec0 = 0$). Nur dann ist sie ein UVR; sonst ein
\emph{affiner} (verschobener) Unterraum.

\begin{fehler}
A1f hast du mit „weiß ich nicht“ gelassen. Antwort: \textbf{falsch} -- eine Hyperebene
ist nur für $n_0=0$ ein Untervektorraum (Open-Book S.\,2 „Untervektorraum“).
\end{fehler}

%==================================================================
\section{SVD richtig aufstellen}

\textbf{Intuition.} Jede Matrix wirkt als: \emph{drehen} ($V\transp$) $\to$
\emph{strecken} ($\Sigma$) $\to$ \emph{drehen} ($U$). $A=U\Sigma V\transp$.

\begin{merke}
\textbf{Dimensionen sind entscheidend:} für $A\in\mathbb{R}^{m\times n}$ ist
$U\in\mathbb{R}^{m\times m}$, $\Sigma\in\mathbb{R}^{m\times n}$ (\emph{gleiche Form wie}
$A$!), $V\in\mathbb{R}^{n\times n}$. $U,V$ sind \emph{orthogonal}; ihre Spalten sind
\textbf{orthonormale Eigenvektoren} -- nicht die Eigenwerte!
\end{merke}

\textbf{Ablauf (Open-Book S.\,1 „SVD“):}
\begin{enumerate}[noitemsep]
  \item $A\transp A$ bilden (symmetrisch, $n\times n$). Eigenwerte $\lambda_i\ge0$,
        normierte Eigenvektoren $\to$ \textbf{Spalten von $V$}.
  \item $\sigma_i=\sqrt{\lambda_i}$, absteigend $\to$ Diagonale von $\Sigma$ (Rest mit
        Nullen auf $m\times n$ auffüllen).
  \item $\vec u_i = \tfrac{1}{\sigma_i}A\vec v_i$ $\to$ Spalten von $U$ (fehlende zu ONB
        ergänzen).
\end{enumerate}

\begin{bsp}
$A=\begin{psmallmatrix}2&0\\0&0\\0&5\end{psmallmatrix}$.\;
$A\transp A=\begin{psmallmatrix}4&0\\0&25\end{psmallmatrix}$,
Eigenwerte $25,4$, also $\sigma_1=5,\ \sigma_2=2$.
Eigenvektoren (absteigend): $\vec v_1=(0,1)\transp$, $\vec v_2=(1,0)\transp$, d.h.
$V=\begin{psmallmatrix}0&1\\1&0\end{psmallmatrix}$.\;
$\vec u_1=\tfrac15A\vec v_1=(0,0,1)\transp$, $\vec u_2=\tfrac12A\vec v_2=(1,0,0)\transp$,
ergänzt $\vec u_3=(0,1,0)\transp$.
\[
U=\begin{psmallmatrix}0&1&0\\0&0&1\\1&0&0\end{psmallmatrix},\quad
\Sigma=\begin{psmallmatrix}5&0\\0&2\\0&0\end{psmallmatrix}\ (3\times2!),\quad
V=\begin{psmallmatrix}0&1\\1&0\end{psmallmatrix}.
\]
\end{bsp}

\begin{fehler}
In A10 hast du $V=\begin{psmallmatrix}4&0\\0&25\end{psmallmatrix}$ gesetzt (das sind die
\emph{Eigenwerte}, nicht die Eigenvektoren!) und $\Sigma$ quadratisch gemacht. Deshalb
ging die Probe schief. Richtig: $V=$ orthonormale Eigenvektoren von $A\transp A$, und
$\Sigma$ hat \textbf{dieselbe Form wie} $A$ ($3\times2$).
\end{fehler}

%==================================================================
\section{Cholesky: die Formel und der häufige Fehler}

\textbf{Intuition.} $A=LL\transp$ ist die „Wurzel“ einer symmetrisch-PD Matrix: ein
einziges Dreieck $L$ statt $L$ und $U$.

\begin{merke}
Open-Book S.\,1 „$LL\transp$“, Reihenfolge $\ell_{11},\ell_{21},\ell_{31},\ell_{22},\dots$:
\[
\ell_{jj}=\sqrt{\,a_{jj}-\textstyle\sum_{k<j}\ell_{jk}^2\,},\qquad
\ell_{ij}=\frac{a_{ij}-\sum_{k<j}\ell_{ik}\ell_{jk}}{\boxed{\ell_{jj}}}\quad(i>j).
\]
\textbf{Geteilt wird durch $\ell_{jj}$ -- NICHT durch $a_{jj}$!}
\end{merke}

\begin{bsp}
$A=\begin{psmallmatrix}4&2&2\\2&5&5\\2&5&14\end{psmallmatrix}$:\;
$\ell_{11}=\sqrt4=2$;\;
$\ell_{21}=\tfrac{2}{\ell_{11}}=\tfrac22=1$;\;
$\ell_{31}=\tfrac{2}{2}=1$;\;
$\ell_{22}=\sqrt{5-1^2}=2$;\;
$\ell_{32}=\tfrac{5-\ell_{31}\ell_{21}}{\ell_{22}}=\tfrac{5-1}{2}=2$;\;
$\ell_{33}=\sqrt{14-1-4}=3$.
\[
L=\begin{psmallmatrix}2&0&0\\1&2&0\\1&2&3\end{psmallmatrix}.
\]
\end{bsp}

\begin{fehler}
In A6b hast du $\ell_{21}=\tfrac{2}{4}=\tfrac12$ gerechnet -- du hast durch $a_{11}=4$
geteilt statt durch $\ell_{11}=\sqrt4=2$. Dadurch lief alles Weitere schief. Immer durch
$\ell_{jj}$ teilen (das ist schon die Wurzel).
\end{fehler}

%==================================================================
\section{Matrixinverse berechnen}

Das hat dir in A5c, A8c und A11 gefehlt -- hier die drei Standardwege.

\subsection{2$\times$2 -- auswendig (Open-Book S.\,2)}
\[
\begin{pmatrix}a&b\\c&d\end{pmatrix}^{-1}
=\frac{1}{ad-bc}\begin{pmatrix}d&-b\\-c&a\end{pmatrix}.
\]
\begin{bsp}(A11) $A\transp A=\begin{psmallmatrix}3&6\\6&14\end{psmallmatrix}$,
$\det=3\cdot14-6\cdot6=6$, also
$(A\transp A)^{-1}=\tfrac16\begin{psmallmatrix}14&-6\\-6&3\end{psmallmatrix}$.
Damit $A^{+}=(A\transp A)^{-1}A\transp
=\begin{psmallmatrix}\tfrac43&\tfrac13&-\tfrac23\\[2pt]-\tfrac12&0&\tfrac12\end{psmallmatrix}$,
und $A^{+}A=E$.\end{bsp}

\subsection{$n\times n$ -- Gauß-Jordan $(A\mid E)\to(E\mid A^{-1})$}
Schreibe $A$ neben die Einheitsmatrix und forme \emph{beide} gleichzeitig um, bis links
$E$ steht; rechts steht dann $A^{-1}$.
\begin{bsp}(A5) Für $A=\begin{psmallmatrix}2&1&1\\4&3&3\\8&7&9\end{psmallmatrix}$ ergibt das
\[
A^{-1}=\begin{psmallmatrix}\tfrac32&-\tfrac12&0\\[2pt]-3&\tfrac52&-\tfrac12\\[2pt]1&-\tfrac32&\tfrac12\end{psmallmatrix}.
\]
Probe: 1.\ Zeile von $A$ mal 1.\ Spalte von $A^{-1}$: $2\cdot\tfrac32+1\cdot(-3)+1\cdot1=1$.~\checkmark
\end{bsp}

\subsection{Über die LU-Zerlegung (das war A5c!)}
\textbf{Idee:} Die $i$-te \emph{Spalte} von $A^{-1}$ ist die Lösung von
$A\vec x_i=\vec e_i$ (denn $A\,A^{-1}=E$, und $E$ hat als Spalten die $\vec e_i$).
Hast du $A=LU$ schon, löse das für jede Spalte in \emph{zwei} Schritten -- ohne neue
Elimination:
\[
  L\vec y=\vec e_i\ (\text{vorwärts}),\qquad U\vec x_i=\vec y\ (\text{rückwärts}).
\]
\begin{bsp}(genau die Matrix aus A5)
$A=\begin{psmallmatrix}2&1&1\\4&3&3\\8&7&9\end{psmallmatrix}$ mit
$L=\begin{psmallmatrix}1&0&0\\2&1&0\\4&3&1\end{psmallmatrix}$,
$U=\begin{psmallmatrix}2&1&1\\0&1&1\\0&0&2\end{psmallmatrix}$.\\[2pt]
\textbf{Spalte 1}, $\vec e_1=(1,0,0)\transp$:\;
$L\vec y=\vec e_1$: $y_1=1$, $2y_1+y_2=0\Rightarrow y_2=-2$,
$4y_1+3y_2+y_3=0\Rightarrow y_3=2$, also $\vec y=(1,-2,2)\transp$.\;
$U\vec x=\vec y$ (rückwärts): $2x_3=2\Rightarrow x_3=1$, $x_2+x_3=-2\Rightarrow x_2=-3$,
$2x_1+x_2+x_3=1\Rightarrow x_1=\tfrac32$. $\Rightarrow$ \textbf{1.\ Spalte} $=(\tfrac32,-3,1)\transp$.\\[2pt]
Genauso mit $\vec e_2,\vec e_3$. Zusammengesetzt:
\[
A^{-1}=\begin{psmallmatrix}\tfrac32&-\tfrac12&0\\[2pt]-3&\tfrac52&-\tfrac12\\[2pt]1&-\tfrac32&\tfrac12\end{psmallmatrix}.
\]
\end{bsp}
\textbf{Merke:} Das ist \emph{derselbe} Vorwärts/Rückwärts-Trick wie beim LGS-Lösen -- nur
dreimal (einmal pro Einheitsvektor). Bei symmetrischem $A=Q\Lambda Q\transp$ geht es noch
schneller: $A^{-1}=Q\Lambda^{-1}Q\transp$ (Kehrwerte auf der Diagonale).

%==================================================================
\section{Abstand Punkt -- Ebene}

\begin{merke}
Ebene $E:\ \vec n\transp\vec x + n_0 = 0$. Für einen Punkt $\vec p$ setze
$w=\vec n\transp\vec p + n_0$. Dann:
\[
d=\frac{|\,\vec n\transp\vec p + n_0\,|}{\norm{\vec n}},\qquad
\text{Seite}=\operatorname{sign}(w),\qquad w=0:\ \text{auf der Ebene.}
\]
\end{merke}

\begin{bsp}
$E:\ 2x_1-x_2+2x_3=3$, also $\vec n=(2,-1,2)\transp$, $n_0=-3$, $\norm{\vec n}=3$.
Punkt $\vec p=(1,2,3)\transp$:\;
$w=(2-2+6)-3=3$, also $d=\tfrac{|3|}{3}=1$.
\end{bsp}

\begin{fehler}
In A13c hast du $d=\tfrac{2p_1-p_2+2p_3}{\norm{\vec n}}$ geschrieben -- der
\textbf{Konstantenterm $n_0=-3$ fehlt}, und du hast nicht zu Ende gerechnet. Richtig:
$|\,\vec n\transp\vec p + n_0\,|/\norm{\vec n}=1$.
\end{fehler}

%==================================================================
\section{Sauber Gauß-eliminieren (Vorzeichen!)}

In A4 und A5 ist dir mehrfach derselbe Fehler passiert: beim Schritt
\textbf{Zielzeile $-$ Faktor$\cdot$Pivotzeile} hast du dich im Vorzeichen vertan
(z.\,B.\ $5-2\cdot2$ als $-1$ statt $1$). Das pflanzt sich in $L$, $U$, $R$, $\det$ fort.

\begin{merke}
$\text{Faktor}=\dfrac{\text{zu eliminierender Eintrag}}{\text{Pivot}}$, dann
\textbf{neue Zeile} $=$ \textbf{alte Zielzeile} $-$ \textbf{Faktor}$\cdot$\textbf{Pivotzeile}.
Die Faktoren stehen \emph{positiv} in $L$ (Open-Book S.\,1 „LU“).
Mach immer die \textbf{Probe} $LU=A$ bzw.\ $CR=A$!
\end{merke}

\begin{bsp}(A5, korrekt)
$\begin{psmallmatrix}2&1&1\\4&3&3\\8&7&9\end{psmallmatrix}$:
$Z_2-2Z_1=(0,1,1)$, $Z_3-4Z_1=(0,3,5)$, dann $Z_3-3Z_2=(0,0,2)$.
\[
L=\begin{psmallmatrix}1&0&0\\2&1&0\\4&3&1\end{psmallmatrix},\quad
U=\begin{psmallmatrix}2&1&1\\0&1&1\\0&0&2\end{psmallmatrix},\quad
\det A = 2\cdot1\cdot2 = 4\ (\text{nicht } -4!).
\]
Du hattest $\ell_{31}=2$ (statt $4$) und Vorzeichen in $U$ falsch.
\end{bsp}

%==================================================================
\section{Finite Differenzen -- nochmal von vorn}

\textbf{Intuition.} Wir können $-u''=f$ nicht exakt am Computer lösen, also ersetzen wir
die Ableitung durch einen \emph{Differenzenquotienten} an wenigen Gitterpunkten und
erhalten ein \textbf{lineares Gleichungssystem} $A\vec u=\vec b$.

\textbf{Baustein (Open-Book S.\,2 „Finite Differenzen“):} zweite Ableitung
\[
u''(x_i)\approx\frac{u_{i-1}-2u_i+u_{i+1}}{h^2}.
\]
Einsetzen in $-u''=f$ ergibt pro innerem Punkt:
$-u_{i-1}+2u_i-u_{i+1}=h^2 f_i$, mit Randwerten $u_0=u_{n+1}=0$.

\begin{bsp}
$-u''=2$, $u(0)=u(1)=0$, $h=\tfrac14$ (Punkte $0{,}25;0{,}5;0{,}75$). $h^2=\tfrac1{16}$,
also $h^2f_i=\tfrac{2}{16}=\tfrac18$. Multipliziert man $-u_{i-1}+2u_i-u_{i+1}=\tfrac18$
mit $16$:
\[
32u_i-16u_{i-1}-16u_{i+1}=2,\quad
A=\begin{psmallmatrix}32&-16&0\\-16&32&-16\\0&-16&32\end{psmallmatrix},\
\vec b=\begin{psmallmatrix}2\\2\\2\end{psmallmatrix}.
\]
$A$ ist \textbf{tridiagonal, symmetrisch, positiv definit, dünn besetzt}.
Lösung (Symmetrie $u_1=u_3$): $u_1=u_3=\tfrac{3}{16}$, $u_2=\tfrac14$.
\end{bsp}

%==================================================================
\section{Lineare Regression -- die Idee}

\textbf{Intuition.} Wir legen die „bestpassende“ Gerade durch Punkte: die Gerade, die die
Summe der quadrierten Abstände minimiert. Das führt \emph{nicht} auf ein lösbares
$X\vec c=\vec y$ (zu viele Gleichungen), sondern auf die \textbf{Normalengleichung}.

\begin{merke}
$X$ = Designmatrix (1.\ Spalte Einsen, dann $x$-Werte), $\vec y$ = Messwerte. Dann
\[
X\transp X\,\vec c = X\transp\vec y\quad\Longrightarrow\quad
\vec c=(X\transp X)^{-1}X\transp\vec y .
\]
$X\transp X$ ist symmetrisch (und bei unabhängigen Spalten PD) -- daher mit
Cholesky/QR lösbar. ($X\transp X$ ist genau das, was die Pseudoinverse $A^{+}$ macht!)
\end{merke}

\begin{bsp}
Punkte $(1,1),(2,2),(3,2)$:
$X=\begin{psmallmatrix}1&1\\1&2\\1&3\end{psmallmatrix}$,
$\vec y=(1,2,2)\transp$,
$X\transp X=\begin{psmallmatrix}3&6\\6&14\end{psmallmatrix}$,
$X\transp\vec y=(5,11)\transp$.
\[
\vec c=\tfrac16\begin{psmallmatrix}14&-6\\-6&3\end{psmallmatrix}\begin{psmallmatrix}5\\11\end{psmallmatrix}
=\begin{psmallmatrix}\tfrac23\\[2pt]\tfrac12\end{psmallmatrix}
\ \Rightarrow\ y=\tfrac23+\tfrac12 x.
\]
\end{bsp}

%==================================================================
\section{Winkel zwischen Vektoren (das war A3d)}

\textbf{Intuition.} Das Skalarprodukt misst, „wie sehr zwei Vektoren in dieselbe Richtung
zeigen“. Aus ihm bekommt man direkt den Winkel.

\begin{merke}
\[
  \cos\alpha = \frac{\inner{\vec u}{\vec v}}{\norm{\vec u}\,\norm{\vec v}},\qquad
  \alpha = \arccos\!\left(\frac{\vec u\transp\vec v}{\norm{\vec u}\,\norm{\vec v}}\right).
\]
Sonderfall: $\inner{\vec u}{\vec v}=0 \Rightarrow \alpha = 90^\circ$ (orthogonal).
\end{merke}

\begin{bsp}
$\vec u=(1,1,0)\transp$, $\vec v=(1,0,1)\transp$:\;
$\inner{\vec u}{\vec v}=1$, $\norm{\vec u}=\norm{\vec v}=\sqrt2$, also
$\cos\alpha=\tfrac{1}{2}\Rightarrow\alpha=60^\circ$.
\end{bsp}

\begin{fehler}
A3d hast du übersprungen (mit der Vermutung, das käme nicht dran). Es ist aber
Standard -- und kommt in der Probeklausur indirekt über das Skalarprodukt vor. Formel
oben, mehr nicht.
\end{fehler}

\vspace{0.6em}{\color{rulecol}\hrule}\vspace{0.4em}
\textit{Tipp: Übe diese Themen mit dem zugehörigen Blatt
\texttt{output/aufgaben/aufgaben\_02.pdf} und teste dich mit
\texttt{output/erklaerung/abfrage\_01.pdf}.}

\end{document}
